% =====================================================================
%  SOLICITUD DE APROBACION ETICA --- Proyecto EAES
%  Dirigida al Vicerrectorado Academico de la UTEQ
%  Compilacion: pdflatex solicitud_comite_etica.tex
% =====================================================================

\documentclass[11pt,a4paper]{article}

\usepackage[T1]{fontenc}
\usepackage[utf8]{inputenc}
\usepackage[spanish,es-noshorthands,es-tabla]{babel}
\usepackage{mathptmx}
\usepackage[scaled=0.90]{helvet}
\usepackage{microtype}
\usepackage{booktabs}
\usepackage{array}
\usepackage{enumitem}
\usepackage{ragged2e}
\usepackage{xcolor}
\usepackage{titlesec}
\usepackage{fancyhdr}
\usepackage{lastpage}

\usepackage[
	a4paper,
	left=2.6cm,
	right=2.6cm,
	top=2.2cm,
	bottom=2.0cm,
	headheight=14pt,
	footskip=1.0cm
]{geometry}

\definecolor{azulinstitucional}{RGB}{20,60,110}
\definecolor{grisclaro}{RGB}{243,245,248}
\definecolor{grismedio}{RGB}{120,130,140}

\titleformat{\section}
	{\normalfont\sffamily\normalsize\bfseries\color{azulinstitucional}}
	{\thesection.}{0.6em}{}
\titlespacing*{\section}{0pt}{9pt}{3pt}

\pagestyle{fancy}
\fancyhf{}
\fancyhead[L]{\footnotesize\sffamily\color{grismedio}Solicitud de aprobaci\'on \'etica --- Proyecto EAES}
\fancyhead[R]{\footnotesize\sffamily\color{grismedio}P\'ag. \thepage\ de \pageref{LastPage}}
\renewcommand{\headrulewidth}{0.4pt}

\setlength{\parskip}{6pt}
\setlength{\parindent}{0pt}
\renewcommand{\arraystretch}{1.35}
\setlist{itemsep=3pt,topsep=3pt,parsep=0pt}

\newcolumntype{L}[1]{>{\RaggedRight\arraybackslash}p{#1}}

\newsavebox{\cajaguardada}
\newenvironment{recuadro}
	{\begin{lrbox}{\cajaguardada}\begin{minipage}{\dimexpr\textwidth-2\fboxsep\relax}\vspace{4pt}\small}
	{\vspace{4pt}\end{minipage}\end{lrbox}\par\vspace{4pt}\noindent\colorbox{grisclaro}{\usebox{\cajaguardada}}\par\vspace{6pt}}

\newcommand{\campo}[1]{\rule{#1}{0.4pt}}

\begin{document}

% =====================================================================
%  OFICIO
% =====================================================================
\begin{center}
	{\sffamily\small\color{grismedio} UNIVERSIDAD T\'ECNICA ESTATAL DE QUEVEDO\par}
	{\sffamily\small\color{grismedio} FACULTAD DE CIENCIAS DE LA COMPUTACI\'ON --- CARRERA DE SOFTWARE\par}
\end{center}

\vspace{6pt}
\hrule
\vspace{10pt}

\begin{flushright}
	\textbf{Oficio Nro.} 001-2026-GGU-FCC-UTEQ

	Quevedo, 12 de agosto de 2026
\end{flushright}

\vspace{6pt}
\textbf{Se\~nor}

\textbf{Dr. Eduardo D\'iaz Ocampo, PhD.}

\textbf{VICERRECTOR ACAD\'EMICO}

\textbf{Universidad T\'ecnica Estatal de Quevedo}

En su despacho.---

\vspace{8pt}
\textbf{Asunto:} Solicitud de revisi\'on y aprobaci\'on \'etica del proyecto de investigaci\'on EA-AES ante el Comit\'e de \'Etica de Investigaci\'on institucional.

\vspace{8pt}
De mi consideraci\'on:

Reciba un cordial saludo. En mi calidad de docente investigador de la Facultad de Ciencias de la Computaci\'on e investigador principal del proyecto que a continuaci\'on se detalla, me dirijo a usted para solicitar, de la manera m\'as comedida, que se sirva disponer su remisi\'on al Comit\'e de \'Etica de Investigaci\'on de la Universidad, a fin de que se realice la revisi\'on correspondiente y, de considerarlo pertinente, se emita el dictamen de aprobaci\'on.

\begin{recuadro}
\textbf{Denominaci\'on del proyecto:} Escritura acad\'emica y Aprendizaje en Educación Superior: estudio multic\'entrico sobre procesos de elaboraci\'on de textos y aprendizaje en educaci\'on superior \textbf{(Proyecto EA-AES)}.

\vspace{4pt}
\textbf{Investigador principal:} Ing. Gleiston Cicer\'on Guerrero Ulloa, PhD.

\vspace{4pt}
\textbf{Unidad ejecutora:} Facultad de Ciencias de la Computaci\'on --- Carrera de Software.

\vspace{4pt}
\textbf{Tipo de estudio:} Experimental, con participantes humanos, aleatorizado, de car\'acter multic\'entrico.

\vspace{4pt}
\textbf{Poblaci\'on participante:} Estudiantes de tercer nivel de diferentes carreras.
\end{recuadro}

El estudio examina c\'omo los estudiantes universitarios elaboran textos acad\'emicos y qu\'e factores inciden en su aprendizaje durante ese proceso. Comprende tres sesiones de escritura en laboratorio supervisado, una prueba posterior de conocimientos y la evaluaci\'on ciega de los textos producidos por parte de revisores externos con experiencia editorial.

Dado que la investigaci\'on involucra participantes humanos, requiere el tratamiento de datos personales y acad\'emicos, y contempla la asignaci\'on aleatoria de los participantes a distintas condiciones de trabajo, considero indispensable contar con el pronunciamiento del Comit\'e de \'Etica antes de iniciar cualquier actividad de reclutamiento o recolecci\'on de informaci\'on.

Se acompa\~na a este oficio el \textbf{memorando t\'ecnico de aspectos \'eticos}, documento elaborado para facilitar la revisi\'on por parte del Comit\'e, en el que se detallan el proceso de consentimiento informado, los riesgos previsibles y sus medidas de mitigaci\'on, el tratamiento y la protecci\'on de los datos, y las razones por las cuales una parte de la informaci\'on sobre las hip\'otesis del estudio se comunicar\'a a los participantes \'unicamente al concluir la investigaci\'on.

\vspace{4pt}
\textbf{Documentaci\'on que se adjunta:}

\begin{enumerate}[label=\textbf{\arabic*.},leftmargin=1.0cm]
	\item Memorando t\'ecnico de aspectos \'eticos.
	\item Protocolo completo de investigaci\'on.
	\item Modelo de consentimiento informado para participantes.
	\item Gu\'ia de escritura acad\'emica que se entregar\'a por igual a todos los participantes.
	\item Compromiso de confidencialidad y destrucci\'on de datos, suscrito por el investigador principal y por el responsable de custodia designado.
	\item Hoja de vida acad\'emica del investigador principal.
\end{enumerate}

\textbf{Declaro expresamente} que no se iniciar\'a el reclutamiento de participantes, la aplicaci\'on de instrumentos ni la solicitud de informaci\'on acad\'emica institucional hasta contar con el dictamen favorable del Comit\'e de \'Etica, y que cualquier modificaci\'on sustantiva del protocolo ser\'a sometida previamente a su conocimiento mediante la enmienda correspondiente.

Por la favorable atenci\'on que se sirva dar al presente, anticipo mi sincero agradecimiento y me suscribo.

\vspace{16pt}
Atentamente,

\vspace{48pt}
\campo{7.6cm}

\textbf{Ing. Gleiston Cicer\'on Guerrero Ulloa, PhD.}

{\small Docente investigador --- Facultad de Ciencias de la Computaci\'on}

{\small Investigador principal del Proyecto EA-AES}

{\small Correo institucional: gguerrero@uteq.edu.ec}

\end{document}
